% gle-BAOAB, one step: the algorithm figure used in README.md.
%
% Scheme (B.1) of Sachs' thesis (Appendix B), assembled from the H-OU splitting of
% section 4.1 with the exact Ornstein-Uhlenbeck step of Eq. (4.7). Written with the
% mass matrix restored: section 4.1 sets M = I throughout.
%
% Build both themes (from this directory):
%
%   pdflatex -jobname=gle-baoab            gle_baoab_algorithm.tex
%   pdflatex -jobname=gle-baoab-dark "\def\DARKTHEME{}% gle-BAOAB, one step: the algorithm figure used in README.md.
%
% Scheme (B.1) of Sachs' thesis (Appendix B), assembled from the H-OU splitting of
% section 4.1 with the exact Ornstein-Uhlenbeck step of Eq. (4.7). Written with the
% mass matrix restored: section 4.1 sets M = I throughout.
%
% Build both themes (from this directory):
%
%   pdflatex -jobname=gle-baoab            gle_baoab_algorithm.tex
%   pdflatex -jobname=gle-baoab-dark "\def\DARKTHEME{}% gle-BAOAB, one step: the algorithm figure used in README.md.
%
% Scheme (B.1) of Sachs' thesis (Appendix B), assembled from the H-OU splitting of
% section 4.1 with the exact Ornstein-Uhlenbeck step of Eq. (4.7). Written with the
% mass matrix restored: section 4.1 sets M = I throughout.
%
% Build both themes (from this directory):
%
%   pdflatex -jobname=gle-baoab            gle_baoab_algorithm.tex
%   pdflatex -jobname=gle-baoab-dark "\def\DARKTHEME{}% gle-BAOAB, one step: the algorithm figure used in README.md.
%
% Scheme (B.1) of Sachs' thesis (Appendix B), assembled from the H-OU splitting of
% section 4.1 with the exact Ornstein-Uhlenbeck step of Eq. (4.7). Written with the
% mass matrix restored: section 4.1 sets M = I throughout.
%
% Build both themes (from this directory):
%
%   pdflatex -jobname=gle-baoab            gle_baoab_algorithm.tex
%   pdflatex -jobname=gle-baoab-dark "\def\DARKTHEME{}\input{gle_baoab_algorithm}"
%   dvisvgm --pdf --no-fonts --exact --output=gle-baoab.svg      gle-baoab.pdf
%   dvisvgm --pdf --no-fonts --exact --output=gle-baoab-dark.svg gle-baoab-dark.pdf
%
% --no-fonts converts glyphs to paths, so the SVG renders identically everywhere:
% GitHub serves it through <img>, where external font loading is blocked.
%
\documentclass[border=10pt]{standalone}

\usepackage[T1]{fontenc}
\usepackage{newtxtext,newtxmath}
\usepackage{inconsolata}
\usepackage{amsmath}
\usepackage{xcolor}
\usepackage[linesnumbered, ruled, vlined]{algorithm2e}

% Palette shared with assets/module_map.py, so the README figures read as one set.
\ifdefined\DARKTHEME
  \definecolor{ink}{HTML}{E6E9EF}
  \definecolor{muted}{HTML}{97A0AF}
  \definecolor{accent}{HTML}{63B6D3}
  \definecolor{bg}{HTML}{0D1117}
\else
  \definecolor{ink}{HTML}{1B2130}
  \definecolor{muted}{HTML}{5F6A7D}
  \definecolor{accent}{HTML}{1D6A87}
  \definecolor{bg}{HTML}{FFFFFF}
\fi

\SetAlgorithmName{gle-BAOAB}{}{}
\renewcommand{\thealgocf}{}                     % no float number: this stands alone
\SetAlgoCaptionSeparator{\,--\,}
\DontPrintSemicolon
% \SetNlSty wraps the number in a ONE-argument macro, so the colour needs its own.
\newcommand{\nlstyle}[1]{\textcolor{muted}{#1}}
\SetNlSty{nlstyle}{}{}                          % line numbers recede
\SetCommentSty{textnormal}
\SetKwInput{KwInit}{Initialise}
\newcommand{\step}[1]{\textcolor{accent}{\textbf{#1}}}
\newcommand{\note}[1]{\textcolor{muted}{\small #1}}

\begin{document}
% \pagecolor, not \colorbox: algorithm2e cannot be scanned as a box argument. The
% standalone page is the cropped figure plus its border, so this paints the whole card.
\pagecolor{bg}
\begin{minipage}{15.4cm}
\color{ink}
\begin{algorithm}[H]
\caption{one step of the quasi-Markovian GLE}
\KwIn{$q^k,\; p^k,\; s^k$;\quad step $\Delta t$;\quad inverse temperature $\beta$;\quad
      $R^k \sim \mathcal{N}(0, I_{n+m})$}
% An input-style line, because the bath is seeded once before the loop rather than
% stepped: \SetKwInput lines are never numbered, whereas \nonl does not hold here.
\KwInit{$s^0 \sim \mathcal{N}(0, \beta^{-1} I_m)$
        \hfill\note{stationary random force}}
\BlankLine
$p^{k+1/2} \leftarrow p^k - \tfrac{\Delta t}{2}\nabla U(q^k)$ \hfill\step{B}\;
$q^{k+1/2} \leftarrow q^k + \tfrac{\Delta t}{2} M^{-1} p^{k+1/2}$ \hfill\step{A}\;
\BlankLine
$\widetilde{\Gamma} \leftarrow \widetilde{\Gamma}(q^{k+1/2})$,\quad
$F \leftarrow e^{-\Delta t \widetilde{\Gamma}}$ \hfill\step{O}\;
\quad take any $S$ with $S S^{\mathsf{T}} = \widetilde{Q} - F \widetilde{Q} F^{\mathsf{T}}$
      \hfill\note{$\widetilde{Q} = I$; pivoted Cholesky}\;
\quad $z \leftarrow \bigl(M^{-1/2} p^{k+1/2},\; s^k\bigr)$\;
\quad $z \leftarrow F z + \beta^{-1/2} S R^k$\;
\quad $\bigl(\hat{p}^{k+1/2},\; s^{k+1}\bigr) \leftarrow \bigl(M^{1/2} z_1,\; z_2\bigr)$\;
\BlankLine
$q^{k+1} \leftarrow q^{k+1/2} + \tfrac{\Delta t}{2} M^{-1} \hat{p}^{k+1/2}$ \hfill\step{A}\;
$p^{k+1} \leftarrow \hat{p}^{k+1/2} - \tfrac{\Delta t}{2}\nabla U(q^{k+1})$ \hfill\step{B}\;
\end{algorithm}
\end{minipage}
\end{document}
"
%   dvisvgm --pdf --no-fonts --exact --output=gle-baoab.svg      gle-baoab.pdf
%   dvisvgm --pdf --no-fonts --exact --output=gle-baoab-dark.svg gle-baoab-dark.pdf
%
% --no-fonts converts glyphs to paths, so the SVG renders identically everywhere:
% GitHub serves it through <img>, where external font loading is blocked.
%
\documentclass[border=10pt]{standalone}

\usepackage[T1]{fontenc}
\usepackage{newtxtext,newtxmath}
\usepackage{inconsolata}
\usepackage{amsmath}
\usepackage{xcolor}
\usepackage[linesnumbered, ruled, vlined]{algorithm2e}

% Palette shared with assets/module_map.py, so the README figures read as one set.
\ifdefined\DARKTHEME
  \definecolor{ink}{HTML}{E6E9EF}
  \definecolor{muted}{HTML}{97A0AF}
  \definecolor{accent}{HTML}{63B6D3}
  \definecolor{bg}{HTML}{0D1117}
\else
  \definecolor{ink}{HTML}{1B2130}
  \definecolor{muted}{HTML}{5F6A7D}
  \definecolor{accent}{HTML}{1D6A87}
  \definecolor{bg}{HTML}{FFFFFF}
\fi

\SetAlgorithmName{gle-BAOAB}{}{}
\renewcommand{\thealgocf}{}                     % no float number: this stands alone
\SetAlgoCaptionSeparator{\,--\,}
\DontPrintSemicolon
% \SetNlSty wraps the number in a ONE-argument macro, so the colour needs its own.
\newcommand{\nlstyle}[1]{\textcolor{muted}{#1}}
\SetNlSty{nlstyle}{}{}                          % line numbers recede
\SetCommentSty{textnormal}
\SetKwInput{KwInit}{Initialise}
\newcommand{\step}[1]{\textcolor{accent}{\textbf{#1}}}
\newcommand{\note}[1]{\textcolor{muted}{\small #1}}

\begin{document}
% \pagecolor, not \colorbox: algorithm2e cannot be scanned as a box argument. The
% standalone page is the cropped figure plus its border, so this paints the whole card.
\pagecolor{bg}
\begin{minipage}{15.4cm}
\color{ink}
\begin{algorithm}[H]
\caption{one step of the quasi-Markovian GLE}
\KwIn{$q^k,\; p^k,\; s^k$;\quad step $\Delta t$;\quad inverse temperature $\beta$;\quad
      $R^k \sim \mathcal{N}(0, I_{n+m})$}
% An input-style line, because the bath is seeded once before the loop rather than
% stepped: \SetKwInput lines are never numbered, whereas \nonl does not hold here.
\KwInit{$s^0 \sim \mathcal{N}(0, \beta^{-1} I_m)$
        \hfill\note{stationary random force}}
\BlankLine
$p^{k+1/2} \leftarrow p^k - \tfrac{\Delta t}{2}\nabla U(q^k)$ \hfill\step{B}\;
$q^{k+1/2} \leftarrow q^k + \tfrac{\Delta t}{2} M^{-1} p^{k+1/2}$ \hfill\step{A}\;
\BlankLine
$\widetilde{\Gamma} \leftarrow \widetilde{\Gamma}(q^{k+1/2})$,\quad
$F \leftarrow e^{-\Delta t \widetilde{\Gamma}}$ \hfill\step{O}\;
\quad take any $S$ with $S S^{\mathsf{T}} = \widetilde{Q} - F \widetilde{Q} F^{\mathsf{T}}$
      \hfill\note{$\widetilde{Q} = I$; pivoted Cholesky}\;
\quad $z \leftarrow \bigl(M^{-1/2} p^{k+1/2},\; s^k\bigr)$\;
\quad $z \leftarrow F z + \beta^{-1/2} S R^k$\;
\quad $\bigl(\hat{p}^{k+1/2},\; s^{k+1}\bigr) \leftarrow \bigl(M^{1/2} z_1,\; z_2\bigr)$\;
\BlankLine
$q^{k+1} \leftarrow q^{k+1/2} + \tfrac{\Delta t}{2} M^{-1} \hat{p}^{k+1/2}$ \hfill\step{A}\;
$p^{k+1} \leftarrow \hat{p}^{k+1/2} - \tfrac{\Delta t}{2}\nabla U(q^{k+1})$ \hfill\step{B}\;
\end{algorithm}
\end{minipage}
\end{document}
"
%   dvisvgm --pdf --no-fonts --exact --output=gle-baoab.svg      gle-baoab.pdf
%   dvisvgm --pdf --no-fonts --exact --output=gle-baoab-dark.svg gle-baoab-dark.pdf
%
% --no-fonts converts glyphs to paths, so the SVG renders identically everywhere:
% GitHub serves it through <img>, where external font loading is blocked.
%
\documentclass[border=10pt]{standalone}

\usepackage[T1]{fontenc}
\usepackage{newtxtext,newtxmath}
\usepackage{inconsolata}
\usepackage{amsmath}
\usepackage{xcolor}
\usepackage[linesnumbered, ruled, vlined]{algorithm2e}

% Palette shared with assets/module_map.py, so the README figures read as one set.
\ifdefined\DARKTHEME
  \definecolor{ink}{HTML}{E6E9EF}
  \definecolor{muted}{HTML}{97A0AF}
  \definecolor{accent}{HTML}{63B6D3}
  \definecolor{bg}{HTML}{0D1117}
\else
  \definecolor{ink}{HTML}{1B2130}
  \definecolor{muted}{HTML}{5F6A7D}
  \definecolor{accent}{HTML}{1D6A87}
  \definecolor{bg}{HTML}{FFFFFF}
\fi

\SetAlgorithmName{gle-BAOAB}{}{}
\renewcommand{\thealgocf}{}                     % no float number: this stands alone
\SetAlgoCaptionSeparator{\,--\,}
\DontPrintSemicolon
% \SetNlSty wraps the number in a ONE-argument macro, so the colour needs its own.
\newcommand{\nlstyle}[1]{\textcolor{muted}{#1}}
\SetNlSty{nlstyle}{}{}                          % line numbers recede
\SetCommentSty{textnormal}
\SetKwInput{KwInit}{Initialise}
\newcommand{\step}[1]{\textcolor{accent}{\textbf{#1}}}
\newcommand{\note}[1]{\textcolor{muted}{\small #1}}

\begin{document}
% \pagecolor, not \colorbox: algorithm2e cannot be scanned as a box argument. The
% standalone page is the cropped figure plus its border, so this paints the whole card.
\pagecolor{bg}
\begin{minipage}{15.4cm}
\color{ink}
\begin{algorithm}[H]
\caption{one step of the quasi-Markovian GLE}
\KwIn{$q^k,\; p^k,\; s^k$;\quad step $\Delta t$;\quad inverse temperature $\beta$;\quad
      $R^k \sim \mathcal{N}(0, I_{n+m})$}
% An input-style line, because the bath is seeded once before the loop rather than
% stepped: \SetKwInput lines are never numbered, whereas \nonl does not hold here.
\KwInit{$s^0 \sim \mathcal{N}(0, \beta^{-1} I_m)$
        \hfill\note{stationary random force}}
\BlankLine
$p^{k+1/2} \leftarrow p^k - \tfrac{\Delta t}{2}\nabla U(q^k)$ \hfill\step{B}\;
$q^{k+1/2} \leftarrow q^k + \tfrac{\Delta t}{2} M^{-1} p^{k+1/2}$ \hfill\step{A}\;
\BlankLine
$\widetilde{\Gamma} \leftarrow \widetilde{\Gamma}(q^{k+1/2})$,\quad
$F \leftarrow e^{-\Delta t \widetilde{\Gamma}}$ \hfill\step{O}\;
\quad take any $S$ with $S S^{\mathsf{T}} = \widetilde{Q} - F \widetilde{Q} F^{\mathsf{T}}$
      \hfill\note{$\widetilde{Q} = I$; pivoted Cholesky}\;
\quad $z \leftarrow \bigl(M^{-1/2} p^{k+1/2},\; s^k\bigr)$\;
\quad $z \leftarrow F z + \beta^{-1/2} S R^k$\;
\quad $\bigl(\hat{p}^{k+1/2},\; s^{k+1}\bigr) \leftarrow \bigl(M^{1/2} z_1,\; z_2\bigr)$\;
\BlankLine
$q^{k+1} \leftarrow q^{k+1/2} + \tfrac{\Delta t}{2} M^{-1} \hat{p}^{k+1/2}$ \hfill\step{A}\;
$p^{k+1} \leftarrow \hat{p}^{k+1/2} - \tfrac{\Delta t}{2}\nabla U(q^{k+1})$ \hfill\step{B}\;
\end{algorithm}
\end{minipage}
\end{document}
"
%   dvisvgm --pdf --no-fonts --exact --output=gle-baoab.svg      gle-baoab.pdf
%   dvisvgm --pdf --no-fonts --exact --output=gle-baoab-dark.svg gle-baoab-dark.pdf
%
% --no-fonts converts glyphs to paths, so the SVG renders identically everywhere:
% GitHub serves it through <img>, where external font loading is blocked.
%
\documentclass[border=10pt]{standalone}

\usepackage[T1]{fontenc}
\usepackage{newtxtext,newtxmath}
\usepackage{inconsolata}
\usepackage{amsmath}
\usepackage{xcolor}
\usepackage[linesnumbered, ruled, vlined]{algorithm2e}

% Palette shared with assets/module_map.py, so the README figures read as one set.
\ifdefined\DARKTHEME
  \definecolor{ink}{HTML}{E6E9EF}
  \definecolor{muted}{HTML}{97A0AF}
  \definecolor{accent}{HTML}{63B6D3}
  \definecolor{bg}{HTML}{0D1117}
\else
  \definecolor{ink}{HTML}{1B2130}
  \definecolor{muted}{HTML}{5F6A7D}
  \definecolor{accent}{HTML}{1D6A87}
  \definecolor{bg}{HTML}{FFFFFF}
\fi

\SetAlgorithmName{gle-BAOAB}{}{}
\renewcommand{\thealgocf}{}                     % no float number: this stands alone
\SetAlgoCaptionSeparator{\,--\,}
\DontPrintSemicolon
% \SetNlSty wraps the number in a ONE-argument macro, so the colour needs its own.
\newcommand{\nlstyle}[1]{\textcolor{muted}{#1}}
\SetNlSty{nlstyle}{}{}                          % line numbers recede
\SetCommentSty{textnormal}
\SetKwInput{KwInit}{Initialise}
\newcommand{\step}[1]{\textcolor{accent}{\textbf{#1}}}
\newcommand{\note}[1]{\textcolor{muted}{\small #1}}

\begin{document}
% \pagecolor, not \colorbox: algorithm2e cannot be scanned as a box argument. The
% standalone page is the cropped figure plus its border, so this paints the whole card.
\pagecolor{bg}
\begin{minipage}{15.4cm}
\color{ink}
\begin{algorithm}[H]
\caption{one step of the quasi-Markovian GLE}
\KwIn{$q^k,\; p^k,\; s^k$;\quad step $\Delta t$;\quad inverse temperature $\beta$;\quad
      $R^k \sim \mathcal{N}(0, I_{n+m})$}
% An input-style line, because the bath is seeded once before the loop rather than
% stepped: \SetKwInput lines are never numbered, whereas \nonl does not hold here.
\KwInit{$s^0 \sim \mathcal{N}(0, \beta^{-1} I_m)$
        \hfill\note{stationary random force}}
\BlankLine
$p^{k+1/2} \leftarrow p^k - \tfrac{\Delta t}{2}\nabla U(q^k)$ \hfill\step{B}\;
$q^{k+1/2} \leftarrow q^k + \tfrac{\Delta t}{2} M^{-1} p^{k+1/2}$ \hfill\step{A}\;
\BlankLine
$\widetilde{\Gamma} \leftarrow \widetilde{\Gamma}(q^{k+1/2})$,\quad
$F \leftarrow e^{-\Delta t \widetilde{\Gamma}}$ \hfill\step{O}\;
\quad take any $S$ with $S S^{\mathsf{T}} = \widetilde{Q} - F \widetilde{Q} F^{\mathsf{T}}$
      \hfill\note{$\widetilde{Q} = I$; pivoted Cholesky}\;
\quad $z \leftarrow \bigl(M^{-1/2} p^{k+1/2},\; s^k\bigr)$\;
\quad $z \leftarrow F z + \beta^{-1/2} S R^k$\;
\quad $\bigl(\hat{p}^{k+1/2},\; s^{k+1}\bigr) \leftarrow \bigl(M^{1/2} z_1,\; z_2\bigr)$\;
\BlankLine
$q^{k+1} \leftarrow q^{k+1/2} + \tfrac{\Delta t}{2} M^{-1} \hat{p}^{k+1/2}$ \hfill\step{A}\;
$p^{k+1} \leftarrow \hat{p}^{k+1/2} - \tfrac{\Delta t}{2}\nabla U(q^{k+1})$ \hfill\step{B}\;
\end{algorithm}
\end{minipage}
\end{document}
